\section{Формирование соглашений}
\subsection{Entente Cordiale}
После второй Англо-бурской войны закончившейся в 1902 году стало понятно, что Англия отягощена своими дальними колониями, и если в них массово начнутся волнения Лондон с этим может не справиться и сколлапсировать. Из-за чего Англия решила сменить свою позицию нейтралитета в Европе на поиск союза. В это время Франция находилась длительное время в изоляции, к тому же у нее были напряженные отношения с Германией. Великобритания также опасалась роста влияния Германии, поэтому империи нашли взаимовыгодным избавиться от противоречий и прийти к сотрудничеству в 1904 они заключили серию соглашений, так называемое "Сердечное согласие" (Entente Cordiale): решили свои противоречия в Египте и Марокко и в других колониях. 

\subsection{Русско-Французский альянс}
В 1890-x была заключена серия соглашения между Францией и Россией, обязавшая империи помочь союзнице в случае войны одной из них с Германией.

\subsection{Англо-Русское соглашение}
В 1907 году между Россией и Англией был заключено соглашение, которое закрывало противоречия между странами в Средней Азии: север Персии отошел России, юг - Британии, центр-нейтральная зона; Тибет под властью Китая, Афганистан сфера влияния Лондона.
Договор был заключен из-за того, что стало понятно: у России больше нет ресурсов на конфликт в регионе (слабость показала Русско-Японская война), к тому же был виден опасный рост силы Германии. Страны нашли более выгодным объединиться против общего врага.

